\documentclass[10pt,a4paper]{amsart}
\usepackage{amsmath,amssymb,amsthm,graphicx}
% Retained minimal preamble for the Stacks Project chapter sites-cohomology.tex (Cohomology on Sites).
% Verbatim excerpt of the chapter preamble: Xy-pic setup (native preamble lines 14-18) and the native
% theorem declarations with the mathematical macros (native preamble lines 154-217). Copied mechanically;
% nothing added, removed or redefined. The wrapper class and the amsmath/amssymb/amsthm/graphicx base
% packages are supplied by the assembler, so the class/externaldocument part of the native preamble is omitted.
% For commutative diagrams we use Xy-pic
\usepackage[all]{xy}

% We use 2cell for 2-commutative diagrams.
\xyoption{2cell}
% Theorem environments.
%
\theoremstyle{plain}
\newtheorem{theorem}[subsection]{Theorem}
\newtheorem{proposition}[subsection]{Proposition}
\newtheorem{lemma}[subsection]{Lemma}

\theoremstyle{definition}
\newtheorem{definition}[subsection]{Definition}
\newtheorem{example}[subsection]{Example}
\newtheorem{exercise}[subsection]{Exercise}
\newtheorem{situation}[subsection]{Situation}

\theoremstyle{remark}
\newtheorem{remark}[subsection]{Remark}
\newtheorem{remarks}[subsection]{Remarks}

\numberwithin{equation}{subsection}

% Macros
%
\def\lim{\mathop{\mathrm{lim}}\nolimits}
\def\colim{\mathop{\mathrm{colim}}\nolimits}
\def\Spec{\mathop{\mathrm{Spec}}}
\def\Hom{\mathop{\mathrm{Hom}}\nolimits}
\def\Ext{\mathop{\mathrm{Ext}}\nolimits}
\def\SheafHom{\mathop{\mathcal{H}\!\mathit{om}}\nolimits}
\def\SheafExt{\mathop{\mathcal{E}\!\mathit{xt}}\nolimits}
\def\Sch{\mathit{Sch}}
\def\Mor{\mathop{\mathrm{Mor}}\nolimits}
\def\Ob{\mathop{\mathrm{Ob}}\nolimits}
\def\Sh{\mathop{\mathit{Sh}}\nolimits}
\def\NL{\mathop{N\!L}\nolimits}
\def\CH{\mathop{\mathrm{CH}}\nolimits}
\def\proetale{{pro\text{-}\acute{e}tale}}
\def\etale{{\acute{e}tale}}
\def\QCoh{\mathit{QCoh}}
\def\Ker{\mathop{\mathrm{Ker}}}
\def\Im{\mathop{\mathrm{Im}}}
\def\Coker{\mathop{\mathrm{Coker}}}
\def\Coim{\mathop{\mathrm{Coim}}}

% Boxtimes
%
\DeclareMathSymbol{\boxtimes}{\mathbin}{AMSa}{"02}

%
% Macros for moduli stacks/spaces
%
\def\QCohstack{\mathcal{QC}\!\mathit{oh}}
\def\Cohstack{\mathcal{C}\!\mathit{oh}}
\def\Spacesstack{\mathcal{S}\!\mathit{paces}}
\def\Quotfunctor{\mathrm{Quot}}
\def\Hilbfunctor{\mathrm{Hilb}}
\def\Curvesstack{\mathcal{C}\!\mathit{urves}}
\def\Polarizedstack{\mathcal{P}\!\mathit{olarized}}
\def\Complexesstack{\mathcal{C}\!\mathit{omplexes}}
% \Pic is the operator that assigns to X its picard group, usage \Pic(X)
% \Picardstack_{X/B} denotes the Picard stack of X over B
% \Picardfunctor_{X/B} denotes the Picard functor of X over B
\def\Pic{\mathop{\mathrm{Pic}}\nolimits}
\def\Picardstack{\mathcal{P}\!\mathit{ic}}
\def\Picardfunctor{\mathrm{Pic}}
\def\Deformationcategory{\mathcal{D}\!\mathit{ef}}

% --- appended 2026-10-12 (Chief): the project own definitions of the reference/slogan/history
% display environments, copied mechanically from the shared Stacks preamble lines 8-12
% (\usepackage{verbatim} plus the three \newenvironment lines, which the chapter preamble keeps
% verbatim from the project preamble). Nothing added, removed or redefined: these environments
% carry no mathematical content and the project itself typesets them as comments. Without these
% three lines a counted statement whose own text uses them fails with Environment slogan undefined.
\usepackage{verbatim}
\newenvironment{reference}{\comment}{\endcomment}
\newenvironment{slogan}{\comment}{\endcomment}
\newenvironment{history}{\comment}{\endcomment}
\usepackage{cleveref}
\crefname{originalreference}{original reference}{original references}
\Crefname{originalreference}{Original reference}{Original references}
\makeatletter
\newlabel{lemma-cone}{{08C5}{source}{}{sourcepointer-1}{}}
\newlabel{lemma-cone@cref}{{[originalreference][0][]08C5}{[1][0][]source}{}{}{}}
\makeatother
\crefname{originalreference}{original reference}{original references}
\Crefname{originalreference}{Original reference}{Original references}
\makeatletter
\newlabel{lemma-cone-pseudo-coherent}{{064R}{source}{}{sourcepointer-2}{}}
\newlabel{lemma-cone-pseudo-coherent@cref}{{[originalreference][0][]064R}{[1][0][]source}{}{}{}}
\makeatother
\crefname{originalreference}{original reference}{original references}
\Crefname{originalreference}{Original reference}{Original references}
\makeatletter
\newlabel{lemma-cone-tor-amplitude}{{0655}{source}{}{sourcepointer-3}{}}
\newlabel{lemma-cone-tor-amplitude@cref}{{[originalreference][0][]0655}{[1][0][]source}{}{}{}}
\makeatother
\crefname{originalreference}{original reference}{original references}
\Crefname{originalreference}{Original reference}{Original references}
\makeatletter
\newlabel{lemma-local-actual}{{08C9}{source}{}{sourcepointer-4}{}}
\newlabel{lemma-local-actual@cref}{{[originalreference][0][]08C9}{[1][0][]source}{}{}{}}
\makeatother
\crefname{originalreference}{original reference}{original references}
\Crefname{originalreference}{Original reference}{Original references}
\makeatletter
\newlabel{lemma-perfect}{{030Z}{source}{}{sourcepointer-5}{}}
\newlabel{lemma-perfect@cref}{{[originalreference][0][]030Z}{[1][0][]source}{}{}{}}
\makeatother
\title{In a distinguished triangle, if two out of three objects are perfect then so is the third}
\author{The Stacks Project Authors}
\date{stacks-project:sites-cohomology.tex:local-snapshot-20261011; GNU FDL 1.2 or later}
\begin{document}
\maketitle
\noindent\textit{Editorial note.} All delivered mathematical text is the author's own, copied from the pinned Stacks Project chapter file sites-cohomology.tex; no statement or proof step was rewritten, reordered or removed.\par\medskip
\noindent\textit{Editorial note.} The counted proof is the author's complete argument: both proofs of the author are delivered: the first combining the characterisation with the triangle lemmas, and the second (which the author labels a sketch) giving the explicit local construction through strictly perfect representatives and cones; the construction contains no omitted step beyond the covering reductions it cites.\par\medskip
\noindent\textit{Editorial note.} The label lemma-two-out-of-three-perfect names the counted statement, and the proof environment immediately following it in the source is delivered verbatim as its proof.\par\medskip
\noindent\textit{Editorial note.} The counted statement is carried by the author's own lemma environment in the chapter source.\par\medskip
\noindent\textit{Editorial note.} This card belongs to the priority queue of the D-72 twin ruling (the sites-side statement is the designated counting side, now counted here for the first time). The statement and proof are copied from the pinned chapter file; the five cited inputs are kept as external pointers with their canonical Stacks tags.\par\medskip
\noindent\textit{Allowed standard prerequisites (editorial).}\begin{itemize}
\item The equivalence between perfection and pseudo-coherence with locally finite tor dimension. Same chapter, the lemma characterising perfect objects as pseudo-coherent objects of locally finite tor dimension. Supplies the framework of the first proof.
\item Pseudo-coherence and tor amplitude in triangles. Same chapter, the lemma on pseudo-coherence in a distinguished triangle with one degree of slack, and the corresponding lemma for tor amplitude. Supplies the two ingredients combined with the characterisation in the first proof.
\item Local representability of maps and strict perfectness of cones. Same chapter, the lemma allowing a map in the derived category to be represented by a map of complexes after a covering, and the lemma that the cone on a morphism of strictly perfect complexes is strictly perfect. Supplies the construction in the second proof.
\end{itemize}\medskip
\ifdefined\hypertarget\hypertarget{sourcepointer-1}{}\fi
\noindent\textit{Source reference.} Original reference 08C5: sources/stacks-sites-cohomology/pinned-original-source.tex, source lines 11735--11735. Prerequisite pointer invoked at proof line(s) 12882; it names a standard fact of the same chapter that the counted argument uses, and every step of the counted proof built on it is retained verbatim in the delivered body.\par\medskip
\ifdefined\hypertarget\hypertarget{sourcepointer-2}{}\fi
\noindent\textit{Source reference.} Original reference 064R: sources/stacks-sites-cohomology/pinned-original-source.tex, source lines 12177--12177. Prerequisite pointer invoked at proof line(s) 12871; it names a standard fact of the same chapter that the counted argument uses, and every step of the counted proof built on it is retained verbatim in the delivered body.\par\medskip
\ifdefined\hypertarget\hypertarget{sourcepointer-3}{}\fi
\noindent\textit{Source reference.} Original reference 0655: sources/stacks-sites-cohomology/pinned-original-source.tex, source lines 12566--12566. Prerequisite pointer invoked at proof line(s) 12872; it names a standard fact of the same chapter that the counted argument uses, and every step of the counted proof built on it is retained verbatim in the delivered body.\par\medskip
\ifdefined\hypertarget\hypertarget{sourcepointer-4}{}\fi
\noindent\textit{Source reference.} Original reference 08C9: sources/stacks-sites-cohomology/pinned-original-source.tex, source lines 11881--11881. Prerequisite pointer invoked at proof line(s) 12880; it names a standard fact of the same chapter that the counted argument uses, and every step of the counted proof built on it is retained verbatim in the delivered body.\par\medskip
\ifdefined\hypertarget\hypertarget{sourcepointer-5}{}\fi
\noindent\textit{Source reference.} Original reference 030Z: sources/stacks-sites-cohomology/pinned-original-source.tex, source lines 12818--12818. Prerequisite pointer invoked at proof line(s) 12871; it names a standard fact of the same chapter that the counted argument uses, and every step of the counted proof built on it is retained verbatim in the delivered body.\par\medskip
\begin{lemma}
\label{lemma-two-out-of-three-perfect}
Let $(\mathcal{C}, \mathcal{O})$ be a ringed site. Let $(K, L, M, f, g, h)$
be a distinguished triangle in $D(\mathcal{O})$. If two out of three of
$K, L, M$ are perfect then the third is also perfect.
\end{lemma}

\begin{proof}
First proof: Combine
Lemmas \ref{lemma-perfect}, \ref{lemma-cone-pseudo-coherent}, and
\ref{lemma-cone-tor-amplitude}.
Second proof (sketch): Say $K$ and $L$ are perfect. Let $U$ be an object
of $\mathcal{C}$. After replacing
$U$ by the members of a covering we may assume that $K|_U$ and $L|_U$
are represented by strictly perfect complexes $\mathcal{K}^\bullet$
and $\mathcal{L}^\bullet$. After replacing $U$ by the members
of a covering we may assume the map $K|_U \to L|_U$ is given by
a map of complexes $\alpha : \mathcal{K}^\bullet \to \mathcal{L}^\bullet$,
see Lemma \ref{lemma-local-actual}.
Then $M|_U$ is isomorphic to the cone of $\alpha$ which is strictly
perfect by Lemma \ref{lemma-cone}.
\end{proof}

\end{document}
